\documentclass[a4paper]{article}

\usepackage[fontset=fandol]{ctex}
\usepackage{amsmath, amssymb}
\usepackage{graphicx}
\usepackage[margin=2.45cm]{geometry}
\usepackage[most]{tcolorbox}
\usepackage{hyperref}
\usepackage{booktabs}
\usepackage{float}
\usepackage{array}
\usepackage{xcolor}
\usepackage{enumitem}
\usepackage{caption}

\setlength{\parindent}{2em}
\setlength{\parskip}{0.35em}
\setlength{\emergencystretch}{3em}
\linespread{1.06}
\sloppy
\raggedbottom
\vfuzz=4pt

\newcolumntype{L}[1]{>{\raggedright\arraybackslash}p{#1}}
\newcolumntype{C}[1]{>{\centering\arraybackslash}p{#1}}

\DeclareMathOperator{\KL}{KL}
\DeclareMathOperator{\softmax}{softmax}
\newcommand{\E}{\mathbb{E}}
\newcommand{\student}{p_{\theta_s}}
\newcommand{\teacher}{q_{\theta_t}}

\hypersetup{
    colorlinks=true,
    linkcolor=blue!60!black,
    urlcolor=blue!60!black
}

\setlist[itemize]{leftmargin=2.2em,itemsep=0.22em,topsep=0.25em}
\setlist[enumerate]{leftmargin=2.4em,itemsep=0.22em,topsep=0.25em}
\setlist[description]{leftmargin=3.0em,itemsep=0.18em,topsep=0.25em}

\newtcolorbox{knowledgebox}[1]{
    enhanced, colback=blue!5!white, colframe=blue!75!black, colbacktitle=blue!75!black,
    coltitle=white, fonttitle=\bfseries, title={#1},
    attach boxed title to top left={yshift=-2mm, xshift=2mm},
    boxrule=1pt, sharp corners
}

\newtcolorbox{importantbox}[1]{
    enhanced, colback=yellow!10!white, colframe=yellow!80!black, colbacktitle=yellow!80!black,
    coltitle=black, fonttitle=\bfseries, title={#1}, sharp corners
}

\newtcolorbox{warningbox}[1]{
    enhanced, colback=red!5!white, colframe=red!75!black, colbacktitle=red!75!black,
    coltitle=white, fonttitle=\bfseries, title={#1}, sharp corners
}

\newcommand{\notetitle}{Network Compression（二）：多视角压缩神经网络}
\newcommand{\notesubtitle}{Knowledge Distillation, Quantization, Efficient Architecture, and Dynamic Computation}
\newcommand{\noteauthors}{基于李宏毅老师《机器学习 2025》课程内容整理}
\newcommand{\notedate}{\today}
\newcommand{\videochannel}{李宏毅深度学习课堂（Bilibili）}
\newcommand{\videoduration}{约 56 分 46 秒}
\newcommand{\videourl}{https://www.bilibili.com/video/BV1TAtwzTE1S?p=52}

\newcommand{\framefig}[4]{%
\begin{figure}[H]
    \centering
    \includegraphics[width=#1\textwidth]{#2}
    \caption{#3\protect\footnotemark}
\end{figure}
\footnotetext{视频画面时间区间：#4。}
}

\begin{document}
\pagestyle{empty}

\begin{center}
    \vspace*{1.0cm}
    {\Huge\bfseries \notetitle\par}
    \vspace{0.35cm}
    {\LARGE \notesubtitle\par}
    \vspace{0.75cm}
    \includegraphics[width=0.62\textwidth]{video.jpg}\par
    \vspace{0.75cm}
    {\large \noteauthors\par}
    {\large \notedate\par}
    \vspace{0.75cm}
    \begin{tcolorbox}[width=0.86\textwidth,center,sharp corners,
                      colback=black!3!white,colframe=black!50!white]
        \centering
        \textbf{视频信息}\\[3pt]
        频道：\videochannel \\
        时长：\videoduration \\
        链接：\href{\videourl}{\videourl}
    \end{tcolorbox}
\end{center}

\newpage
\tableofcontents
\newpage
\pagestyle{plain}

% =====================================================================
\section{从 Pruning 到多路线压缩}
% =====================================================================

上一讲聚焦 network pruning：先训练一个较大的模型，再移除低重要性的权重、神经元或结构块，并通过 fine-tuning 恢复性能。本讲继续讨论 network compression，但不再只看 pruning，而是从多个面向压缩神经网络：让小模型模仿大模型、用更少 bit 表示参数、直接设计高效架构，以及让模型按实际需求动态改变计算量。

课程最后的总结页把这一讲的范围列成五个主题：
\begin{itemize}
    \item \textbf{Network Pruning：} 从结构中删除冗余部分。
    \item \textbf{Knowledge Distillation：} 把大模型或 ensemble 的行为转移给小模型。
    \item \textbf{Parameter Quantization：} 用较少 bit 表示权重或 activation。
    \item \textbf{Architecture Design：} 重新设计参数效率更高、计算更省的网络模块。
    \item \textbf{Dynamic Computation：} 不同输入、设备或状态使用不同计算量。
\end{itemize}

\framefig{0.82}{figures/fig14_concluding_remarks.jpg}
{本讲最后回顾的 compression 路线：pruning、distillation、quantization、architecture design 与 dynamic computation。}
{00:56:15--00:56:30}

这些路线不是互斥的。实际系统中常见组合是：先选一个更高效的 backbone，再用大模型 distillation 训练它，随后做 quantization，必要时再剪枝或加入 early exit。也就是说，压缩不是某个单一技巧，而是围绕部署约束组织的一组工程与学习方法。

\begin{importantbox}{压缩的目标不是“越小越好”}
    模型压缩真正优化的是多目标 trade-off：准确率、模型大小、延迟、吞吐、内存访问、能耗、硬件支持、训练成本与维护复杂度。一个方法若只减少参数量，却没有减少实际 wall-clock latency，在部署上未必有价值。
\end{importantbox}

\subsection{本章小结}

Network compression 可以从结构、数值表示、学习目标、网络模块和运行路径几个层面同时入手。P51 讲的是 pruning，本讲补上其它常见方向，让压缩从“删掉一些东西”扩展成“重新组织模型的知识、表示和计算”。

% =====================================================================
\section{Knowledge Distillation：把大模型的行为教给小模型}
% =====================================================================

Knowledge Distillation 的基本设定是：先有一个能力较强但较大的 teacher network，再训练一个较小的 student network，让 student 不只学习真实标签，还学习 teacher 对输入的完整输出分布。teacher 可以是单个大模型，也可以是多个模型 ensemble 的平均输出。

\framefig{0.86}{figures/fig01_kd_teacher_student_target.jpg}
{Knowledge distillation 的核心流程：teacher net 对输入图片给出软目标，例如 ``1'' 的概率为 0.7、``7'' 为 0.2、``9'' 为 0.1，student net 通过 cross-entropy 或 KL divergence 去匹配这个分布。}
{00:02:15--00:02:30}

传统监督学习通常只给 one-hot 标签。比如某张图的标签是数字 ``1''，训练目标就是 $y_1=1$，其它类别为 $0$。但 teacher 的输出分布更细：它可能认为这张图确实最像 ``1''，但也有一点像 ``7'' 或 ``9''。这种“次高概率”的信息常被称为 \textbf{dark knowledge}：它不在硬标签中显式出现，却表达了类别之间的相似性和 teacher 学到的判别边界。

若 teacher 在 temperature $T$ 下输出分布 $\teacher^T(y\mid x)$，student 在同一 temperature 下输出 $\student^T(y\mid x)$，蒸馏损失常写成：
\[
    L_{\mathrm{KD}}
    =
    T^2 \KL\!\left(\teacher^T(\cdot\mid x)\,\Vert\,\student^T(\cdot\mid x)\right).
\]
实际训练中也常把它和真实标签的监督损失相加：
\[
    L
    =
    \alpha L_{\mathrm{hard}}(y, \student)
    +
    (1-\alpha)L_{\mathrm{KD}},
\]
其中 $L_{\mathrm{hard}}$ 是普通 cross-entropy，$\alpha$ 控制硬标签与 teacher soft label 的权重。

\begin{knowledgebox}{为什么 soft label 比 hard label 更丰富}
    如果 hard label 只告诉 student “这是 1”，student 无法知道它和 7、9 的相似程度。teacher 的 soft label 则告诉 student：这张图主要像 1，但局部笔画也可能使它有一点像 7 或 9。对小模型来说，这种相似度结构能减少盲目探索，也能把大模型的 inductive bias 转移过来。
\end{knowledgebox}

\subsection{从单个 teacher 到 ensemble teacher}

teacher 不一定是一张网络。课程图中也给出 ensemble 版本：把很多网络的输出取平均，再拿平均后的分布训练 student。ensemble 往往比单模型更稳健，因为不同模型会犯不同错误，平均后可以降低方差。问题是 ensemble 部署成本高：推理时要跑很多模型，参数量和延迟都会成倍增加。

\framefig{0.86}{figures/fig02_kd_ensemble_teacher.jpg}
{Ensemble distillation：先平均多个 teacher network 的输出分布，再让一个小 student 学习这个平均分布，从而把 ensemble 的预测行为压缩到单个模型中。}
{00:07:15--00:07:30}

Distillation 因此可以理解为一种“行为压缩”：不是直接把 teacher 的参数复制给 student，而是把 teacher 的 input-output function 通过数据和目标分布转移出去。它和 pruning 的精神也有相通处：都先依赖一个较大的模型提供能力或知识，再把可部署模型变小。

\subsection{Temperature：让 softmax 变得更软}

softmax 对 logit 的尺度很敏感。如果 logit 差距很大，普通 softmax 会把分布压成接近 one-hot，次要类别的概率几乎为零，dark knowledge 就被隐藏掉了。temperature 的做法是先把 logit 除以 $T$：
\[
    p_i^{(T)}
    =
    \frac{\exp(z_i/T)}
    {\sum_j \exp(z_j/T)}.
\]
当 $T>1$ 时，类别之间的 logit 差距被缩小，softmax 输出更平滑；当 $T$ 很大时，原本近乎 one-hot 的分布会显露出更多相对关系。

\framefig{0.84}{figures/fig03_kd_temperature_softmax.jpg}
{Temperature softmax 的效果：原本 $z=(100,10,1)$ 会给出接近 one-hot 的分布；除以 $T=100$ 后，logit 变成 $(1,0.1,0.01)$，输出约为 $(0.56,0.23,0.21)$，次要类别信息被保留下来。}
{00:11:00--00:11:15}

训练 student 时通常用较高 temperature 计算 teacher 和 student 的 soft distribution；推理时 student 仍可用普通 $T=1$ 的输出。这里要注意，temperature 不是为了让模型“更不确定”，而是为了让训练信号包含类别间关系。若 teacher 本身很差，或者 soft target 被过度平滑，student 也可能学到噪声。

\begin{warningbox}{Distillation 不是免费的性能保证}
    小模型容量有限，无法完整复制大模型。蒸馏是否有效取决于 student 架构、teacher 质量、训练数据覆盖、temperature、硬标签与软标签损失权重等因素。如果 student 太小，soft label 可能只能提供有限帮助；如果 teacher 的偏差很强，student 也会继承这些偏差。
\end{warningbox}

\subsection{本章小结}

Knowledge distillation 用 teacher 的输出分布训练 student，把大模型或 ensemble 的行为压缩到小模型中。soft label 的价值在于保留类别相似性，temperature 则让这些相似性更容易被 student 看到。它压缩的不是参数本身，而是模型学到的函数行为。

% =====================================================================
\section{Parameter Quantization：用更少 bit 表示参数}
% =====================================================================

Parameter Quantization 关注的是数值表示。神经网络中大量权重通常用 32-bit 浮点数存储，但很多任务并不需要每个参数都保留这么高的精度。如果能用 16-bit、8-bit，甚至更少 bit 表示权重，就能减少模型文件大小和内存带宽需求；若硬件支持低精度计算，还可能降低延迟和能耗。

课程先给出最直接的想法：\textbf{using less bits to represent a value}。例如把连续权重映射到有限 codebook，实际存储时只存 codebook index。若 codebook 有 $K$ 个中心，只需要 $\lceil \log_2 K\rceil$ bit 表示一个权重，再额外保存一个很小的查表。

\framefig{0.86}{figures/fig04_quantization_weight_clustering.jpg}
{Weight clustering 的量化方式：把原始权重聚成少数几类，矩阵中每个位置只保存 cluster index；若只有 4 个 cluster，每个位置只需 2 bit，再用 table 还原代表值。频繁 cluster 还可以用 Huffman encoding 给更短码。}
{00:21:00--00:21:15}

图中把多个实数权重映射成四种颜色，每种颜色对应一个代表值。这样做有两个层次的压缩：
\begin{itemize}
    \item \textbf{聚类压缩：} 参数本身不再保存完整浮点数，只保存 cluster id。
    \item \textbf{熵编码压缩：} 高频 cluster 用更短 bit pattern，低频 cluster 用更长 bit pattern，例如 Huffman encoding。
\end{itemize}

量化可以发生在训练后，也可以发生在训练中。训练后量化（post-training quantization）实现简单，但若精度降得太狠，模型可能无法适应。量化感知训练（quantization-aware training）会在训练时模拟量化误差，让模型提前学会在低精度约束下工作。

\begin{importantbox}{量化要同时看存储、带宽和硬件 kernel}
    参数变成 8-bit 并不自动意味着推理加速。如果计算图中仍频繁反量化回浮点，或者目标硬件没有高效的 int8 kernel，实际速度可能提升有限。量化最稳定的收益通常来自模型文件更小、内存带宽更低；延迟收益需要结合硬件和框架验证。
\end{importantbox}

\subsection{Binary Weights：极端的量化}

Binary weights 是更激进的量化：权重只允许 $+1$ 或 $-1$。这样模型存储极小，乘法也可能简化为符号操作。课程中提到 BinaryConnect、BinaryNet、XNOR-Net 等工作，它们尝试在训练或推理时使用二值权重。

BinaryConnect 的训练思路可以概括为：前向传播和反向传播中用二值权重计算网络行为和梯度方向，但参数更新仍维护一份 real-valued weights。也就是说，二值网络负责“参与计算”，实数权重负责“累计更新”。这可以避免权重只有 $\pm 1$ 时更新空间过于僵硬。

\framefig{0.86}{figures/fig05_binaryconnect_training_mechanism.jpg}
{BinaryConnect 的训练机制示意：network with binary weights 负责计算负梯度，真实值权重根据更新方向移动；训练时保留 real-valued weights，使用时可二值化。}
{00:22:15--00:22:30}

课程也展示了 BinaryConnect 在 MNIST、CIFAR-10、SVHN 上的实验表。表格中 BinaryConnect 的误差率并未灾难性上升，甚至某些设置下接近或优于没有 regularizer 的 baseline。这说明在过参数化网络中，很多任务对权重精度并没有想象中那么敏感。

\framefig{0.86}{figures/fig06_binary_weights_results.jpg}
{Binary weights 的实验结果示意：BinaryConnect 在多个数据集上能达到接近 baseline 的表现，显示出神经网络对低精度参数的容忍度。}
{00:23:30--00:23:45}

不过，二值化是非常强的约束。它更适合特定架构、特定硬件和特定任务；对大型语言模型、多模态模型或高精度回归任务，直接二值化往往会带来明显损失。现代部署中更常见的是 int8、int4、混合精度量化，以及对不同层使用不同 bit 的方法。

\subsection{本章小结}

Parameter quantization 通过降低数值精度压缩模型。Weight clustering 把连续权重映射到 codebook，binary weights 则把这种思想推到极端。量化的关键不只是“bit 数变少”，还包括是否能训练稳定、是否能保持精度、是否能被目标硬件真正利用。

% =====================================================================
\section{Architecture Design：直接设计更省的网络}
% =====================================================================

Pruning 和 quantization 都是在已有模型上做压缩；architecture design 则从一开始就设计更高效的模块。课程用 convolution 为例说明：标准卷积的参数量和计算量往往很大，而适当拆解卷积结构，可以在保留表达力的同时显著减少参数。

\subsection{标准卷积的参数量}

设输入 feature map 有 $I$ 个 channel，输出有 $O$ 个 channel，卷积核大小为 $k\times k$。标准 convolution 中，每个 output channel 都需要一个覆盖所有 input channel 的卷积核，因此参数量为
\[
    k \times k \times I \times O.
\]
课程图中例子为 $k=3$、$I=2$、$O=4$，所以参数量是
\[
    3\times 3\times 2\times 4 = 72.
\]

\framefig{0.84}{figures/fig07_standard_cnn_parameters.jpg}
{标准 CNN 卷积参数量：每个 output channel 都连接所有 input channel，因此参数量为 $k^2IO$；图中 $3\times3\times2\times4=72$。}
{00:30:30--00:30:45}

标准卷积的问题不是它“不好”，而是它同时混合了两类计算：空间维度上的局部特征抽取，以及 channel 之间的信息混合。Depthwise separable convolution 的想法就是把这两件事拆开做。

\subsection{Depthwise Separable Convolution}

Depthwise separable convolution 分为两步：
\begin{enumerate}
    \item \textbf{Depthwise convolution：} 每个 input channel 各自使用一个 $k\times k$ filter，只在本 channel 内做空间卷积，参数量为 $k^2I$。
    \item \textbf{Pointwise convolution：} 用 $1\times1$ convolution 混合 channel，产生 $O$ 个 output channel，参数量为 $IO$。
\end{enumerate}
总参数量从标准卷积的 $k^2IO$ 变为
\[
    k^2I + IO.
\]
两者比例为
\[
    \frac{k^2I+IO}{k^2IO}
    =
    \frac{1}{O}+\frac{1}{k^2}.
\]
当 $k=3$ 且 $O$ 较大时，比例约为 $1/9$ 加一个很小的项，参数和计算都可明显下降。

\framefig{0.84}{figures/fig08_depthwise_separable_parameters.jpg}
{Depthwise separable convolution 的参数量对比：标准卷积为 $k^2IO$，拆成 depthwise 与 pointwise 后为 $k^2I+IO$，比例为 $\frac{1}{O}+\frac{1}{k^2}$。}
{00:37:15--00:37:30}

Depthwise separable convolution 是 MobileNet、Xception 等高效 CNN 架构的重要基础。它牺牲的是“每个空间 filter 直接跨所有 channel 混合”的自由度，换来更低参数量和更少计算。现代高效网络往往会在这种模块基础上继续加入 inverted bottleneck、channel shuffle、squeeze-and-excitation 等设计。

\begin{knowledgebox}{模块设计也是一种压缩}
    如果一个架构从训练开始就使用更高效的模块，它不需要先训练一个笨重模型再压缩。它把压缩目标写进了模型的 inductive bias：哪些交互值得保留，哪些交互可以分解，哪些计算可以共享。
\end{knowledgebox}

\subsection{Low-Rank Approximation}

Low-rank approximation 是另一类常见结构压缩。若全连接层权重矩阵 $W\in\mathbb{R}^{M\times N}$ 近似低秩，可以把它分解为两个较小矩阵：
\[
    W \approx U V,\qquad
    U\in\mathbb{R}^{M\times K},\quad
    V\in\mathbb{R}^{K\times N},\quad
    K<\min(M,N).
\]
原本参数量是 $MN$，分解后是 $MK+KN$。只要 $K$ 足够小，参数就会明显减少。从网络结构角度看，这相当于在原本一层 linear mapping 中插入一个 bottleneck hidden dimension。

\framefig{0.82}{figures/fig09_low_rank_factorization.jpg}
{Low-rank approximation：把大矩阵 $W$ 近似为 $U$ 与 $V$ 的乘积，中间 rank $K$ 小于 $M,N$，从而用两个小矩阵替代一个大矩阵。}
{00:41:00--00:41:15}

Low-rank 思想也可以用于 convolution：把一个大卷积分解成多个较小卷积，或者把 channel mixing 与 spatial filtering 拆开。与 pruning 相比，low-rank 分解得到的结构通常更规整，更容易被常规矩阵乘法或卷积 kernel 利用。

\subsection{本章小结}

Architecture design 从模型结构本身降低计算量。标准卷积的 $k^2IO$ 参数可以通过 depthwise separable convolution 降到 $k^2I+IO$；大矩阵也可以用 low-rank factorization 近似。它们共同说明：压缩不只是在训练后删参数，也可以在设计阶段改变网络的计算分解方式。

% =====================================================================
\section{Dynamic Computation：按需求改变计算量}
% =====================================================================

前面几类方法通常得到一个固定大小、固定计算图的模型。Dynamic computation 则进一步问：是否每次推理都必须运行同样多的 layer、channel 或 block？现实中，不同设备、不同电量状态、不同输入难度，可能需要不同计算量。

\framefig{0.84}{figures/fig10_dynamic_computation_motivation.jpg}
{Dynamic computation 的动机：模型可以根据设备类型、电量状态或任务需求调整计算量，而不是对所有情境都运行同一套固定计算。}
{00:47:15--00:47:30}

一个朴素办法是准备多个模型：手机低电量用小模型，云端或高电量用大模型。但这样会增加训练、存储和维护成本。Dynamic computation 希望用一个模型覆盖多种计算预算，让它在运行时选择深度、宽度或路径。

\subsection{Dynamic Depth：中途出口}

Dynamic depth 在网络中加入多个 early exit。浅层输出 $y'$、中层输出 $y''$、最终输出 $y$ 都可以接受监督。训练时把每个出口的误差加起来：
\[
    L = e_1 + e_2 + \cdots + e_L.
\]
这样浅层出口也学会做预测。推理时，如果资源有限或中途预测已经足够可信，就可以提前停止；如果需要更高精度，则继续往后跑。

\framefig{0.84}{figures/fig11_dynamic_depth_exits.jpg}
{Dynamic depth：在中间层接出 extra layers 形成多个出口，低预算时可以使用较浅输出，高预算时继续通过完整网络；训练损失可写为各出口误差之和。}
{00:50:45--00:51:00}

Dynamic depth 的关键难点是出口判定。若模型太早退出，可能因为浅层特征不足而犯错；若总是不退出，又失去节省计算的意义。常见做法包括根据 confidence、entropy、margin 或 learned controller 判断是否继续。

\subsection{Dynamic Width：可切换宽度}

Dynamic width 让同一个网络支持不同宽度。宽模式使用更多神经元或 channel，窄模式只使用其中一部分。课程图中给出 Slimmable Neural Networks：同一网络可以对应多种宽度，每一种宽度都有自己的输出和误差项：
\[
    L = e_1 + e_2 + e_3.
\]
训练时同时约束多个宽度，使子网络也能独立工作。

\framefig{0.84}{figures/fig12_dynamic_width_slimmable.jpg}
{Dynamic width：同一网络共享参数，但可以启用不同数量的神经元或 channel，形成从宽到窄的多个子网络。}
{00:52:15--00:52:30}

这种方法比维护多套模型更省存储，但训练更复杂。宽网络和窄网络共享权重，彼此的梯度可能互相干扰；batch normalization、activation range 和 calibration 也可能随宽度变化。工程实现中通常需要额外技巧保证不同宽度都稳定。

\subsection{Computation Based on Sample Difficulty}

Dynamic computation 还可以根据样本难度调整。简单样本经过浅层就能正确判断，可以早停；困难样本则继续经过更多层。课程图中用两张图说明：一张简单小猫图在浅层已经预测为 cat，可以 stop；另一张较难图浅层误判为 Mexican rolls，因此不能停，必须继续往深层处理。

\framefig{0.86}{figures/fig13_sample_difficulty_early_exit.jpg}
{按样本难度分配计算：简单样本浅层输出已足够可信时早停；困难样本若浅层判断错误或不确定，则继续运行后续层。}
{00:55:15--00:55:30}

这类方法的直觉很强：不是所有输入都同样难。对于图像分类，清晰、典型的样本可能很早就被识别；模糊、遮挡、分布外或类别相近的样本则需要更深层特征。课程提到的相关方向包括 SkipNet、Runtime Neural Pruning 与 BlockDrop，它们都在探索如何为不同输入选择不同推理路径。

\begin{warningbox}{Early exit 需要可靠的不确定性估计}
    如果 confidence 本身不可靠，模型可能对错误答案非常自信，从而过早退出。Dynamic computation 在节省计算前，必须保证退出策略经过校准和验证，否则会把一部分错误隐藏在“低成本推理”中。
\end{warningbox}

\subsection{本章小结}

Dynamic computation 让模型在推理时根据设备、资源或样本难度改变计算量。Dynamic depth 改变运行层数，dynamic width 改变启用宽度，sample-dependent routing 则按输入选择路径。它的优势是灵活，但核心风险是控制策略：什么时候少算，什么时候必须多算。

% =====================================================================
\section{总结与延伸}
% =====================================================================

本讲把 network compression 从 pruning 扩展到更完整的技术地图。不同方法压缩的是不同对象：
\begin{itemize}
    \item \textbf{Pruning 压缩结构：} 删除权重、神经元、channel、head 或 block。
    \item \textbf{Distillation 压缩行为：} 用小模型学习大模型或 ensemble 的输出分布。
    \item \textbf{Quantization 压缩表示：} 用更少 bit 存储和计算权重、activation。
    \item \textbf{Architecture design 压缩计算分解：} 用更高效的模块替代昂贵模块。
    \item \textbf{Dynamic computation 压缩每次推理：} 对不同情境使用不同计算预算。
\end{itemize}

这些方法背后的共同问题是：神经网络中哪些部分是“必须保留的能力”，哪些只是训练便利、数值冗余或固定计算图带来的浪费？Pruning 说冗余连接可以删；distillation 说大模型的行为可以转教给小模型；quantization 说参数不必用高精度表示；efficient architecture 说计算可以重新分解；dynamic computation 说每个样本也不必走同样路径。

\begin{importantbox}{部署视角下的推荐思路}
    真实应用中，通常先明确瓶颈：是模型文件太大、内存带宽太高、延迟太高、能耗太高，还是不同设备需要不同预算。然后再选方法：文件大小优先可先看 quantization；延迟优先要看 architecture 与 structured pruning；精度损失难以接受时可加入 distillation；输入难度差异明显时再考虑 early exit 或 dynamic routing。
\end{importantbox}

从学习角度看，compression 也揭示了深度模型的一个重要事实：训练时需要的大容量，和部署时最终需要的容量，不一定相同。大模型可能是一个强 teacher、一个优化过程中的宽搜索空间，或一个可被提炼的知识载体；压缩方法则尝试把真正有用的部分留下来，让模型在现实约束下运行。

\subsection{本章小结}

P52 的核心不是某个单独技巧，而是 compression 的系统视角：参数、结构、数值、知识与运行路径都可以被压缩。后续真正做模型部署时，应把这些方法当作工具箱，根据任务、硬件和误差预算组合使用，而不是只追求某一个压缩率数字。

\end{document}
